\documentclass[11pt, a4paper]{article}

\usepackage[utf8]{inputenc}
\usepackage[T1]{fontenc}
\usepackage[french]{babel}
\usepackage{amsmath, amssymb}
\usepackage{geometry}
\usepackage{graphicx}
\usepackage{tabularx}
\usepackage{booktabs}
\usepackage{xcolor}

% Configuration des marges pour correspondre à l'original
\geometry{
    top=2cm,
    bottom=2cm,
    left=1.5cm,
    right=1.5cm
}

% Commande personnalisée pour les placeholders d'images
\newcommand{\imageplaceholder}[2]{
    \begin{center}
        \fbox{\begin{minipage}{#1}\centering\vspace{1cm}\textit{#2}\vspace{1cm}\end{minipage}}
    \end{center}
}

% Commande pour afficher la correction avec explication
\newcommand{\corr}[1]{
    \vspace{0.3cm}
    \begin{center}
    \fcolorbox{blue}{blue!5}{
    \begin{minipage}{0.95\textwidth}
    \color{blue} \textbf{Correction \& Raisonnement :} \\
    #1
    \end{minipage}}
    \end{center}
    \vspace{0.3cm}
}

\begin{document}

% ==================== EN-TÊTE ====================
\noindent
\begin{tabularx}{\textwidth}{@{}X c r@{}}
\textbf{EP360} & \textbf{Mardi 19 mai 2026} & \textbf{Sans documents} \\
\textbf{Electronique} & \textbf{Examen 2h} & \textbf{Avec calculette} \\
\end{tabularx}
\vspace{0.2cm}
\hrule
\vspace{0.4cm}

% ==================== PROBLÈME I ====================
\begin{center}
    \textbf{Problème I - 10 pts}
\end{center}
\vspace{-0.2cm}
\hrule
\vspace{0.3cm}

\noindent
\begin{minipage}[t]{0.52\textwidth}
a) Donner la représentation générale d'un transistor bipolaire en paramètres hybrides, sous forme matricielle.\\
On considèrera pour la suite $h_{12}=0$ et $h_{22}=0$.

\corr{
\textbf{Étape 1 : Modèle général du quadripôle.} \\
La représentation en paramètres hybrides permet de lier les grandeurs d'entrée ($v_{be}$, $i_b$) aux grandeurs de sortie ($v_{ce}$, $i_c$) d'un transistor[cite: 3]. La matrice générale s'écrit :
$$ \begin{bmatrix} v_{be} \\ i_c \end{bmatrix} = \begin{bmatrix} h_{11} & h_{12} \\ h_{21} & h_{22} \end{bmatrix} \begin{bmatrix} i_b \\ v_{ce} \end{bmatrix} $$ [cite: 4]

\textbf{Étape 2 : Application des hypothèses simplificatrices.} \\
Le sujet nous donne $h_{12}=0$ et $h_{22}=0$. Les équations deviennent donc :
$$ v_{be} = h_{11} \cdot i_b $$
$$ i_c = h_{21} \cdot i_b $$

\textbf{Conclusion :}\\
Le transistor se modélise par un circuit simple : une résistance d'entrée $h_{11}$ et une source de courant commandée $h_{21}i_b$ (souvent notée $\beta i_b$)[cite: 4].
}

b) Donner le schéma équivalent dynamique du montage de la Fig. 1 vu entre les bornes B, C et E.

\corr{
\textbf{Étape 1 : Remplacement par les modèles dynamiques.}\\
Pour établir le schéma équivalent, il faut remplacer chaque transistor ($T_1$ et $T_2$) par le modèle simplifié déterminé à la question a)[cite: 5].

\textbf{Étape 2 : Identification des connexions (Nœuds).}\\
En suivant le schéma de la Fig. 1, on réalise les connexions suivantes :
\begin{itemize}
    \item \textbf{Base globale (B) :} Elle correspond directement à la base du premier transistor ($B_1$)[cite: 6].
    \item \textbf{Liaison interne ($E_1$ - $B_2$) :} L'émetteur de $T_1$ est directement connecté à la base de $T_2$[cite: 7]. 
    \item \textbf{Collecteur global (C) :} Les collecteurs $C_1$ et $C_2$ sont reliés ensemble pour former le collecteur global C[cite: 8].
    \item \textbf{Émetteur global (E) :} Il correspond à l'émetteur du deuxième transistor ($E_2$)[cite: 8].
\end{itemize}
\textit{Le schéma comporte donc $h_{11,T1}$ en série avec $h_{11,T2}$, et deux sources de courant ($h_{21,T1}ib_1$ et $h_{21,T2}ib_2$) injectant en parallèle dans le collecteur commun C[cite: 8].}
}

c) Exprimer $ib_2$ en fonction de $ib_1$ et de $h_{21,T1}$ (on ne négligera pas le terme 1 devant $h_{21}$).

\corr{
\textbf{Étape 1 : Analyse de la liaison $E_1$-$B_2$.}\\
Le courant entrant dans la base $B_2$ ($ib_2$) correspond exactement au courant d'émetteur de $T_1$ ($i_{e1}$)[cite: 9].

\textbf{Étape 2 : Loi des nœuds sur le transistor 1.}\\
La somme des courants entrant dans un transistor est égale au courant sortant par l'émetteur[cite: 10]:
$$ i_{e1} = i_{b1} + i_{c1} $$ [cite: 10]

\textbf{Étape 3 : Substitution et factorisation.}\\
On sait que $i_{c1} = h_{21,T1} \cdot i_{b1}$[cite: 10]. En remplaçant dans l'équation :
$$ ib_2 = i_{b1} + h_{21,T1} \cdot i_{b1} $$ [cite: 11]
$$ \mathbf{ib_2 = (1 + h_{21,T1}) \cdot i_{b1}} $$ [cite: 11]
}

d) En déduire l'expression de l'impédance d'entrée équivalente.

\corr{
\textbf{Étape 1 : Définir l'impédance d'entrée.}\\
L'impédance d'entrée globale est définie par $h_{11,Teq} = \frac{v_{be}}{i_b}$[cite: 12].

\textbf{Étape 2 : Loi des mailles à l'entrée.}\\
La maille d'entrée donne : $v_{be} = v_{be1} + v_{be2}$[cite: 12]. En utilisant les modèles hybrides pour chaque transistor :
$$ v_{be} = h_{11,T1} \cdot i_{b1} + h_{11,T2} \cdot i_{b2} $$ [cite: 13]

\textbf{Étape 3 : Substitution de $ib_2$.}\\
On remplace $ib_2$ par l'expression obtenue en c) :
$$ v_{be} = h_{11,T1} \cdot i_{b1} + h_{11,T2} \cdot (1 + h_{21,T1}) \cdot i_{b1} $$ [cite: 13]

\textbf{Étape 4 : Factorisation.}\\
En factorisant par $i_{b1}$ (sachant que le courant d'entrée global $i_b = i_{b1}$)[cite: 13]:
$$ \mathbf{h_{11,Teq} = h_{11,T1} + h_{11,T2}(1 + h_{21,T1})} $$ [cite: 13]
}

e) En déduire l'expression du gain en courant équivalent.

\corr{
\textbf{Étape 1 : Définir le gain en courant.}\\
Le gain en courant global est $h_{21,Teq} = \frac{i_c}{i_b}$[cite: 14].

\textbf{Étape 2 : Loi des nœuds au collecteur.}\\
Le courant de collecteur total est la somme des deux courants de collecteur[cite: 14]:
$$ i_c = i_{c1} + i_{c2} = h_{21,T1} \cdot i_{b1} + h_{21,T2} \cdot i_{b2} $$ [cite: 14]

\textbf{Étape 3 : Substitution et développement.}\\
On remplace $ib_2$ par son expression :
$$ i_c = h_{21,T1} \cdot i_{b1} + h_{21,T2} \cdot (1 + h_{21,T1}) \cdot i_{b1} $$ [cite: 14]
En factorisant par $i_{b1}$ :
$$ \mathbf{h_{21,Teq} = h_{21,T1} + h_{21,T2} + h_{21,T1} \cdot h_{21,T2}} $$ [cite: 14]
}

f) Dans le cas où les deux transistors sont identiques, simplifier les expressions de : $h_{11,Teq}$ et $h_{21,Teq}$

\corr{
\textbf{Hypothèse :}\\
Les transistors sont identiques, on pose donc $h_{11,T1} = h_{11,T2} = h_{11}$ et $h_{21,T1} = h_{21,T2} = \beta$[cite: 14].

\textbf{Simplification de $h_{11,Teq}$ :}\\
$$ h_{11,Teq} = h_{11} + h_{11}(1 + \beta) = h_{11}(2 + \beta) $$ [cite: 14, 15]
Comme $\beta \gg 1$ en général, on approxime[cite: 15]:
$$ \mathbf{h_{11,Teq} \approx \beta \cdot h_{11}} $$ [cite: 15]

\textbf{Simplification de $h_{21,Teq}$ :}\\
$$ h_{21,Teq} = \beta + \beta + \beta^2 = 2\beta + \beta^2 $$ [cite: 15]
Avec $\beta \gg 1$, le terme au carré est dominant :
$$ \mathbf{h_{21,Teq} \approx \beta^2} $$ [cite: 15]
}
\end{minipage}
\hfill
\begin{minipage}[t]{0.45\textwidth}
g) Indiquer de quel type de montage il s'agit et préciser son principal intérêt.

\corr{
\textbf{Type de montage :}\\
Il s'agit d'un \textbf{Montage Darlington}[cite: 16]. 

\textbf{Intérêt principal :}\\
Ce montage permet de se comporter comme un transistor unique possédant :
\begin{itemize}
    \item Un \textbf{très grand gain en courant} (environ $\beta^2$)[cite: 16].
    \item Une \textbf{très grande impédance d'entrée}[cite: 16].
\end{itemize}
Cela le rend idéal pour des étages de puissance ou d'amplification nécessitant un très faible courant de commande[cite: 16].
}

\imageplaceholder{0.9\textwidth}{[Insérer Figure 1 : Montage Darlington]}
\begin{center}
    {\small Fig. 1}
\end{center}
\end{minipage}

\vspace{0.5cm}

% ==================== PROBLÈME II ====================
\begin{center}
    \textbf{Problème II - 10 pts}
\end{center}
\vspace{-0.2cm}
\hrule
\vspace{0.3cm}

\noindent
\begin{minipage}[t]{0.45\textwidth}
Le filtre ci-dessous est constitué d'un AOP parfait, alimenté en symétrique associé à deux quadripôles $\text{Q}_\text{A}$, $\text{Q}_\text{B}$.
Les quadripôles A et B sont décrits par leur matrice admittance.
\imageplaceholder{0.9\textwidth}{[Insérer Figures 2 et 3 : Schéma AOP et quadripôles]}
\begin{center}
    {\small Fig. 2 \hspace{2cm} Fig. 3}
\end{center}

$$ \text{Y}_\text{A} = \begin{bmatrix} y_{A_{11}} & y_{A_{12}} \\ y_{A_{21}} & y_{A_{22}} \end{bmatrix} \qquad \text{Y}_\text{B} = \begin{bmatrix} y_{B_{11}} & y_{B_{12}} \\ y_{B_{21}} & y_{B_{22}} \end{bmatrix} $$

a/ Exprimer en démontrant H(p) en fonction de certains éléments $y_{x\;ij}$ : $\text{H(p)} = \frac{v_S}{v_E}$ \quad avec (x = A ou B).

\corr{
\textbf{Étape 1 : Propriétés de l'AOP parfait en régime linéaire.}\\
L'AOP possède une boucle de rétroaction sur son entrée inverseuse (via $Q_B$), il fonctionne donc en régime linéaire. Par conséquent, la tension différentielle est nulle : $V^+ = V^-$.
Puisque l'entrée non-inverseuse est à la masse ($V^+ = 0$), on obtient une **masse virtuelle** : $V^- = 0$.
De plus, les courants d'entrée d'un AOP parfait sont nuls : $i^- = 0$.

\textbf{Étape 2 : Équations de transfert des quadripôles.}\\
\begin{itemize}
    \item \textbf{Pour le quadripôle A} (connecté entre $V_e$ et $V^-$) : le courant de sortie $I_{A2}$ s'écrit, d'après la matrice admittance : 
    $$ I_{A2} = y_{A21} V_e + y_{A22} V^- $$
    Comme $V^- = 0$, cela se simplifie en : $I_{A2} = y_{A21} V_e$.
    
    \item \textbf{Pour le quadripôle B} (connecté entre $V^-$ et $V_s$) : le courant d'entrée $I_{B1}$ s'écrit :
    $$ I_{B1} = y_{B11} V^- + y_{B12} V_s $$
    Comme $V^- = 0$, cela se simplifie en : $I_{B1} = y_{B12} V_s$.
\end{itemize}

\textbf{Étape 3 : Loi des nœuds à l'entrée inverseuse.}\\
Le courant sortant de $Q_A$ s'additionne au courant entrant dans $Q_B$, et la somme doit être égale au courant $i^-$ de l'AOP :
$$ I_{A2} + I_{B1} = i^- = 0 $$

\textbf{Étape 4 : Déduction de la fonction de transfert.}\\
En remplaçant les courants par les expressions trouvées à l'étape 2 :
$$ y_{A21} V_e + y_{B12} V_s = 0 $$
$$ y_{B12} V_s = - y_{A21} V_e $$
$$ \mathbf{H(p) = \frac{V_s}{V_e} = -\frac{y_{A21}}{y_{B12}}} $$
}
\end{minipage}
\hfill
\begin{minipage}[t]{0.52\textwidth}
b/ Exprimez H(p) la fonction de transfert du filtre en considérant :\\
\begin{tabular}{ll}
$y_1$ : capacité C & $y_2$ : capacité C \\
$y_3$ : résistance R & $y_4$ : résistance R \\
\end{tabular}

\vspace{0.2cm}
c/ Donner la fonction de transfert sous forme canonique en fonction des composants.

\vspace{0.2cm}
d/ Définissez l'ordre et le type de filtre.

\vspace{0.4cm}
Données\\
\hspace*{0.5cm} $\cdot$ Matrice Admittance -- Quadripôle A
$$ \text{Y}_\text{A} = \frac{1}{y_1+y_2+y_3} \begin{bmatrix} y_1(y_2+y_3) & -y_1 y_2 \\ -y_1 y_2 & y_2(y_1+y_3) \end{bmatrix} $$

\hspace*{0.5cm} $\cdot$ Matrice Admittance -- Quadripôle B
$$ \text{Y}_\text{B} = \begin{bmatrix} \dots & \frac{-y_1 y_2}{y_1+y_2+y_3} - y_4 \\[10pt] \dots & \dots \end{bmatrix} $$

\corr{
\textbf{Questions b/ et c/ :}\\
\textbf{Étape 1 : Identification des paramètres des matrices.}\\
En observant les matrices $Y_A$ et $Y_B$ fournies, on repère les coefficients nécessaires :
$$ y_{A21} = \frac{-y_1 y_2}{y_1+y_2+y_3} \quad \text{et} \quad y_{B12} = \frac{-y_1 y_2}{y_1+y_2+y_3} - y_4 $$

\textbf{Étape 2 : Remplacement dans H(p).}\\
D'après la question a/ :
$$ H(p) = -\frac{\frac{-y_1 y_2}{y_1+y_2+y_3}}{\frac{-y_1 y_2}{y_1+y_2+y_3} - y_4} $$
En multipliant numérateur et dénominateur par $(y_1+y_2+y_3)$ pour éliminer les fractions étagées :
$$ H(p) = \frac{y_1 y_2}{y_1 y_2 + y_4(y_1+y_2+y_3)} $$

\textbf{Étape 3 : Remplacement par les admittances des composants complexes.}\\
L'admittance d'un condensateur est $Cp$ et celle d'une résistance est $1/R$. Ainsi, $y_1 = y_2 = Cp$ et $y_3 = y_4 = \frac{1}{R}$.
$$ H(p) = \frac{(Cp)(Cp)}{(Cp)(Cp) + \frac{1}{R}\left(Cp + Cp + \frac{1}{R}\right)} $$
$$ H(p) = \frac{C^2p^2}{C^2p^2 + \frac{1}{R}\left(2Cp + \frac{1}{R}\right)} = \frac{C^2p^2}{C^2p^2 + \frac{2C}{R}p + \frac{1}{R^2}} $$

\textbf{Étape 4 : Mise sous forme canonique.}\\
Pour identifier une forme canonique standard, on cherche à obtenir "1" comme terme de degré 0 au dénominateur. On multiplie numérateur et dénominateur par $R^2$ :
$$ \mathbf{H(p) = \frac{(RCp)^2}{1 + 2RCp + (RCp)^2}} $$

\vspace{0.2cm}
\textbf{Question d/ :}\\
\textbf{Étape 1 : Identification avec la forme générale.}\\
La fonction de transfert correspond à la forme canonique :
$$ H(p) = \frac{A_0 \cdot \left(\frac{p}{\omega_0}\right)^2}{1 + 2m\left(\frac{p}{\omega_0}\right) + \left(\frac{p}{\omega_0}\right)^2} $$
Avec $\omega_0 = \frac{1}{RC}$, $A_0 = 1$ et $m = 1$.

\textbf{Étape 2 : Conclusion sur le filtre.}\\
\begin{itemize}
    \item \textbf{Ordre :} La puissance maximale de $p$ au dénominateur est 2. C'est un filtre du \textbf{2ème ordre}.
    \item \textbf{Type :} À très basse fréquence ($p \to 0$), $H(p) \to 0$. À très haute fréquence ($p \to \infty$), $H(p) \to \frac{(RCp)^2}{(RCp)^2} = 1$. Il laisse donc passer les hautes fréquences, c'est un filtre \textbf{passe-haut}.
\end{itemize}
}
\end{minipage}

\vspace{0.5cm}

% ==================== PROBLÈME III ====================
\begin{center}
    \textbf{Problème III - 10 pts}
\end{center}
\vspace{-0.2cm}
\hrule
\vspace{0.3cm}

\noindent
\begin{minipage}[t]{0.45\textwidth}
\imageplaceholder{0.95\textwidth}{[Insérer Figure 4 : Filtre universel (AOP multiples)]}
\end{minipage}
\hfill
\begin{minipage}[t]{0.52\textwidth}
Les AOP sont parfaits alimentés symétriques +/- $\text{V}_{\text{cc}}$.
\begin{itemize}
    \item Exprimez la tension $v_2$ en fonction de $V_e$ et $V_1$.
    \item Exprimez la tension $v_s$ en fonction de $V_2$.
    \item Exprimez la tension $v_1$ en fonction de $V_s$.
    \item Exprimez T(p) la fonction de transfert sous forme canonique en exprimant la pulsation propre, le coefficient d'amortissement et le gain statique : $\text{T(p)} = \frac{v_s}{v_e}$
    \item Quel est le type et l'ordre du filtre
    \item Pourquoi tant de composants pour ce filtre ?
\end{itemize}
\end{minipage}

\vspace{0.3cm}

\noindent
\begin{minipage}{\textwidth}
\corr{
\textbf{1. Expression de la tension $v_2$ en fonction de $V_e$ et $V_1$}\\
\textbf{Étape 1 : Analyse du régime de fonctionnement de $A_1$.}\\
L'AOP $A_1$ possède une boucle de rétroaction négative via la résistance $R$, il fonctionne donc en régime linéaire. Comme son entrée non-inverseuse est à la masse ($V^+ = 0$), on en déduit une masse virtuelle sur l'entrée inverseuse : $V^- = 0$. Les courants d'entrée de l'AOP étant nuls, on peut appliquer la loi des nœuds en ce point.\\
\textbf{Étape 2 : Application de la loi des nœuds.}\\
Trois branches de composants arrivent au nœud $V^-$. Comme $V^- = 0$, on somme les courants sortant de ce nœud vers $V_e$, vers $V_1$ et vers $v_2$ :
$$ Cp(0 - V_e) + Cp(0 - V_1) + \frac{0 - v_2}{R} = 0 $$
\textbf{Étape 3 : Résolution algébrique.}\\
$$ -CpV_e - CpV_1 - \frac{v_2}{R} = 0 $$
$$ \frac{v_2}{R} = -Cp(V_e + V_1) \implies \mathbf{v_2 = -RCp(V_e + V_1)} $$

\vspace{0.2cm}
\textbf{2. Expression de la tension $v_s$ en fonction de $V_2$}\\
\textbf{Étape 1 : Analyse du nœud inverseur de $A_2$.}\\
De la même manière, $A_2$ fonctionne en régime linéaire avec $V^+ = 0$, donc son entrée inverseuse est également à un potentiel nul ($V^- = 0$).\\
\textbf{Étape 2 : Détermination de l'admittance de la boucle de rétroaction.}\\
La rétroaction de $A_2$ est composée d'une résistance $R$ en parallèle avec un condensateur $C_3$. L'admittance globale de cette branche parallèle est : $Y_f = \frac{1}{R} + C_3p$.\\
\textbf{Étape 3 : Application de la loi des nœuds au point $V^- = 0$.}\\
Le courant arrivant de $v_2$ à travers le condensateur $C$ est égal au courant repartant vers $v_s$ à travers la branche parallèle :
$$ (v_2 - 0)\cdot Cp + (v_s - 0)\cdot \left(\frac{1}{R} + C_3p\right) = 0 $$
$$ v_s \cdot \left(\frac{1 + RC_3p}{R}\right) = -v_2 \cdot Cp \implies \mathbf{v_s = -v_2 \frac{RCp}{1 + RC_3p}} $$

\vspace{0.2cm}
\textbf{3. Expression de la tension $v_1$ en fonction de $V_s$}\\
\textbf{Étape 1 : Identification du montage $A_3$.}\\
L'AOP $A_3$ est câblé selon une structure classique d'amplificateur inverseur.\\
\textbf{Étape 2 : Équation du montage.}\\
Le signal d'entrée est $v_s$ (appliqué sur la résistance d'entrée $R_1$) et la sortie est $v_1$ (liée par la résistance de rétroaction $R_1$). La formule du gain de l'inverseur donne directement :
$$ v_1 = -\frac{R_1}{R_1} \cdot v_s \implies \mathbf{v_1 = -v_s} $$

\vspace{0.2cm}
\textbf{4. Expression de la fonction de transfert globale $T(p) = \frac{v_s}{v_e}$}\\
\textbf{Étape 1 : Substitution des variables intermédiaires.}\\
Remplaçons d'abord $V_1$ par $-v_s$ (trouvé à la question 3) dans l'expression de $v_2$ (question 1) :
$$ v_2 = -RCp(V_e - v_s) = RCp(v_s - V_e) $$
\textbf{Étape 2 : Injection dans l'équation de la tension de sortie.}\\
Introduisons cette expression de $v_2$ dans la relation de la question 2 :
$$ v_s = - \left[ RCp(v_s - V_e) \right] \cdot \frac{RCp}{1 + RC_3p} = (V_e - v_s) \frac{(RCp)^2}{1 + RC_3p} $$
\textbf{Étape 3 : Isolation de la variable de sortie $v_s$.}\\
Multiplions toute l'équation par $(1 + RC_3p)$ :
$$ v_s(1 + RC_3p) = V_e(RCp)^2 - v_s(RCp)^2 \implies v_s \left[1 + RC_3p + (RCp)^2\right] = V_e(RCp)^2 $$
On obtient la fonction de transfert globale :
$$ \mathbf{T(p) = \frac{v_s}{V_e} = \frac{(RCp)^2}{1 + RC_3p + (RCp)^2}} $$
\textbf{Étape 4 : Identification des paramètres de la forme canonique.}\\
La forme générale standard d'un filtre du second ordre passe-haut est : $T(p) = \frac{T_0 \left(\frac{p}{\omega_0}\right)^2}{1 + 2m\left(\frac{p}{\omega_0}\right) + \left(\frac{p}{\omega_0}\right)^2}$. En posant $\frac{p}{\omega_0} = RCp$, on en déduit par identification :
\begin{itemize}
    \item Pulsation propre : $\mathbf{\omega_0 = \frac{1}{RC}}$
    \item Gain statique : $\mathbf{T_0 = 1}$
    \item Coefficient d'amortissement ($m$) : le terme du premier degré donne $2m\frac{p}{\omega_0} = 2mRCp$. En égalisant avec notre terme $RC_3p$, on obtient $2mRC = RC_3 \implies \mathbf{m = \frac{C_3}{2C}}$.
\end{itemize}

\vspace{0.2cm}
\textbf{5. Type et ordre du filtre}\\
\textbf{Étape 1 : Détermination de l'ordre.}\\
Le dénominateur de la fonction de transfert est un polynôme en $p$ dont le degré maximal est de 2. Il s'agit donc d'un filtre du \textbf{2ème ordre}.\\
\textbf{Étape 2 : Détermination du comportement fréquentiel.}\\
\begin{itemize}
    \item À basse fréquence ($p \to 0$) : Le numérateur tend vers 0, donc $T(p) \to 0$.
    \item À haute fréquence ($p \to \infty$) : Les termes de plus haut degré l'emportent, $T(p) \to \frac{(RCp)^2}{(RCp)^2} = 1$.
\end{itemize}
Le filtre bloque les basses fréquences et laisse passer les hautes fréquences, c'est un filtre \textbf{passe-haut}.

\vspace{0.2cm}
\textbf{6. Intérêt de cette structure (Filtre universel / Variable d'état)}\\
Bien que ce circuit nécessite 3 AOP et de nombreux composants passifs, il offre un avantage industriel majeur par rapport aux structures plus simples (Sallen-Key ou Rauch) : \textbf{l'indépendance des réglages}.\\
En effet, l'équation montre que :
\begin{itemize}
    \item On peut modifier la fréquence de coupure $\omega_0$ en faisant varier la résistance $R$ sans perturber le facteur d'amortissement $m$.
    \item On peut ajuster le facteur d'amortissement $m$ uniquement en modifiant la capacité $C_3$, sans impacter la fréquence de coupure $\omega_0$.
\end{itemize}
Cette décorrélation rend le filtre extrêmement stable, précis et facile à calibrer.
}
\end{minipage}

\vspace{0.5cm}
% ==================== PROBLÈME IV ====================
\begin{center}
    \textbf{Problème IV - 10 pts}
\end{center}
\vspace{-0.2cm}
\hrule
\vspace{0.3cm}

\noindent
\begin{minipage}[t]{0.45\textwidth}
La fonction de transfert en boucle fermée du système présenté à la Fig. 5 s'écrit :
$$ T(j\omega) = \frac{A(j\omega)}{1 - A(j\omega).B(j\omega)} $$

\imageplaceholder{0.8\textwidth}{[Insérer Figure 5 : Schéma-bloc oscillateur]}
\begin{center}
    {\small Fig. 5}
\end{center}
\end{minipage}
\hfill
\begin{minipage}[t]{0.52\textwidth}
On veut réaliser un oscillateur avec ce système qui est constitué :\\
d'un ampli inverseur de gain $A = -K$\\
et d'un réseau de retour ayant comme fonction de transfert :
$$ B(j\omega) = \frac{1}{\left(1+j\frac{\omega}{\omega_0}\right)^3} \qquad \text{avec } \omega_0 = 10^4 \text{ rad/s} $$

1. Tracer le diagramme de Bode (module et phase) de $B(j\omega)$.\\
2. Determiner graphiquement et par calcul 
la pulsation d'oscillation $\omega_{osc}$. Justifier.\\
3. Quel est le gain minimal K permettant l'oscillation ?
\end{minipage}

\vspace{0.3cm}

\noindent
\begin{minipage}{\textwidth}
\corr{
\textbf{1. Établissement du diagramme de Bode de $B(j\omega)$}\\
\textbf{Étape 1 : Analyse du module et du gain en dB.}\\
Le module de la fonction de transfert s'exprime par : 
$$ |B(j\omega)| = \frac{1}{\left(\sqrt{1 + \left(\frac{\omega}{\omega_0}\right)^2}\right)^3} = \left(1 + \left(\frac{\omega}{\omega_0}\right)^2\right)^{-3/2} $$
Le gain en décibels est donc : $G_{dB}(\omega) = 20\log|B(j\omega)| = -30\log\left(1 + \left(\frac{\omega}{\omega_0}\right)^2\right)$.
\begin{itemize}
    \item \textbf{À basse fréquence} ($\omega \ll \omega_0$) : $G_{dB}(\omega) \approx 0 \text{ dB}$ (asymptote horizontale).
    \item \textbf{À haute fréquence} ($\omega \gg \omega_0$) : $G_{dB}(\omega) \approx -60\log\left(\frac{\omega}{\omega_0}\right)$, ce qui donne une droite de pente \textbf{$-60 \text{ dB/décade}$} (comportement d'un système d'ordre 3).
    \item \textbf{À la pulsation propre} ($\omega = \omega_0$) : $G_{dB}(\omega_0) = -30\log(2) \approx -9 \text{ dB}$.
\end{itemize}

\textbf{Étape 2 : Analyse de la phase.}\\
L'argument de $B(j\omega)$ est donné par la somme des arguments (le dénominateur apporte un signe moins) :
$$ \phi(\omega) = \arg(B(j\omega)) = -3 \arctan\left(\frac{\omega}{\omega_0}\right) $$
\begin{itemize}
    \item \textbf{À basse fréquence} ($\omega \to 0$) : $\phi \to -3 \times 0^\circ = 0^\circ$.
    \item \textbf{À la pulsation propre} ($\omega = \omega_0$) : $\phi = -3 \times \arctan(1) = -3 \times 45^\circ = -135^\circ$.
    \item \textbf{À haute fréquence} ($\omega \to \infty$) : $\phi \to -3 \times 90^\circ = -270^\circ$.
\end{itemize}

\vspace{0.3cm}
\textbf{2. Détermination de la pulsation d'oscillation $\omega_{osc}$}\\
\textbf{Étape 1 : Condition théorique d'oscillation (Critère de Barkhausen).}\\
Pour qu'un circuit oscille de façon autonome (sans signal d'entrée), le système en boucle fermée doit avoir un pôle sur l'axe imaginaire, ce qui se traduit par un dénominateur nul : $1 - A\cdot B = 0 \implies A\cdot B(j\omega_{osc}) = 1$.\\
En remplaçant l'amplificateur par son gain $A = -K$, on obtient : 
$$ -K \cdot B(j\omega_{osc}) = 1 \implies B(j\omega_{osc}) = -\frac{1}{K} $$

\textbf{Étape 2 : Extraction de la condition de phase.}\\
Le gain $K$ étant un réel positif, le terme $-\frac{1}{K}$ est un nombre réel purement négatif. Son argument est donc égal à $\pm 180^\circ$ (ou $\pi$ radians)[cite: 45]. On cherche la pulsation $\omega_{osc}$ telle que :
$$ \arg(B(j\omega_{osc})) = -180^\circ $$

\textbf{Étape 3 : Calcul mathématique de $\omega_{osc}$.}\\
Utilisons l'expression de la phase trouvée à la question 1 :
$$ -3 \arctan\left(\frac{\omega_{osc}}{\omega_0}\right) = -180^\circ \implies \arctan\left(\frac{\omega_{osc}}{\omega_0}\right) = 60^\circ \quad \left(\text{ou } \frac{\pi}{3} \text{ rad}\right) $$
En appliquant la fonction tangente des deux côtés :
$$ \frac{\omega_{osc}}{\omega_0} = \tan(60^\circ) = \sqrt{3} \implies \mathbf{\omega_{osc} = \omega_0\sqrt{3}} $$
En effectuant l'application numérique avec $\omega_0 = 10^4 \text{ rad/s}$ :
$$ \mathbf{\omega_{osc} = 10^4 \times \sqrt{3} \approx 17\,320 \text{ rad/s}} $$

\textbf{Étape 4 : Justification graphique.}\\
Sur le diagramme de Bode de phase, la pulsation $\omega_{osc}$ correspond précisément au point d'intersection où la courbe de phase croise la ligne horizontale d'ordonnée $-180^\circ$.

\vspace{0.3cm}
\textbf{3. Calcul du gain minimal $K$ permettant le démarrage de l'oscillation}\\
\textbf{Étape 1 : Application de la condition de module.}\\
D'après le critère de Barkhausen ($A \cdot B = 1$), la condition sur les modules impose : $|A| \cdot |B(j\omega_{osc})| = 1 \implies K \cdot |B(j\omega_{osc})| = 1$. Pour assurer le démarrage des oscillations face aux pertes du circuit, le gain réel doit être supérieur ou égal à cette limite :
$$ K \ge \frac{1}{|B(j\omega_{osc})|} $$

\textbf{Étape 2 : Calcul de la valeur du module de retour à $\omega_{osc}$.}\\
À la pulsation d'oscillation, nous avons établi que $\frac{\omega_{osc}}{\omega_0} = \sqrt{3}$. Remplaçons ce rapport dans la formule du module de $B$ :
$$ |B(j\omega_{osc})| = \frac{1}{\left(\sqrt{1 + (\sqrt{3})^2}\right)^3} = \frac{1}{\left(\sqrt{1 + 3}\right)^3} = \frac{1}{(\sqrt{4})^3} = \frac{1}{2^3} = \frac{1}{8} $$

\textbf{Étape 3 : Déduction du gain limite.}\\
$$ K \ge \frac{1}{1/8} \implies \mathbf{K_{min} = 8} $$
Il faut donc un gain d'amplification d'au moins \textbf{8} pour compenser l'atténuation d'un facteur 8 subie par le signal lors de son passage dans le réseau de retour à la fréquence d'oscillation.
}
\end{minipage}

\vfill
% ==================== PIED DE PAGE ====================
\hrule
\vspace{0.1cm}
\noindent
Allane \hfill Esisar \hfill 1/2

\end{document}
